% ===========================================================================
\section{Related work}
\label{sec:related}

\paragraph{Self-preference and self-recognition.}
LLM judges rate their own outputs above others'
\citep{panickssery2024selfpreference, wataoka2024perplexity, ye2025justice}.
\citet{panickssery2024selfpreference} tie the preference to self-recognition
through fine-tuning, and \citet{ackerman2025inspection} locate a
self-recognition direction in Llama-3-8B-Instruct.
\citet{wataoka2024perplexity} propose that the preference follows
familiarity: judges favour text of low perplexity to them, which their own
text usually is. Whether judges can \emph{name} their own text is less
settled. \citet{bai2025selfrecognition} find exact-author accuracy near chance
for ten models shown one text at a time, with guesses piling onto GPT and
Claude; \citet{davidson2024anonymized} find no consistent self-recognition in
anonymised panels; \citet{zhou2025implicit} find that models recognise their
own text in pairs but not alone; and \citet{laine2024sad} include
self-recognition in a situational-awareness benchmark.
\citet{stamand2026sgtr} show that a quality heuristic, picking the
better text as one's own, confounds many of these measurements, and
\citet{roytburg2026narcissists} and \citet{shi2025positionbias} document how
judge biases shift with format and position. Two studies are closest to
ours. CALM \citep{ye2025justice} scores self-enhancement against the scores
\emph{other} judges give the same answer, a peer reference for quality; we
apply the same idea to naming the author. \citet{barkhordar2026stylenotself}
show, for code, that raw attribution accuracy mostly reflects how readily a
model claims authorship and recommend balanced accuracy and heuristic
baselines. We reach the same diagnosis for prose and three other genres, and
add the peer baseline, the comparison of formats on identical texts, and
likelihood access for open-weight judges.

\paragraph{Where the self is located.}
\citet{khullar2026selfattribution} define \emph{self-attribution bias} as a
monitor rating an action more favourably when it is implicitly framed as its
own; the effect is strongest when the monitor judges an action it produced
earlier in the same conversation and largely vanishes for other models'
actions. That is a judgment shift driven by conversational position, not an
authorship decision, which is why we avoid the term and speak of self-naming.
Our on-policy condition (Section~\ref{sec:onpolicy}) connects the two: it
asks whether a judge that has just written an answer recognises it as its
own.

\paragraph{Attributing text to models.}
Supervised classifiers separate model families with high accuracy
\citep{sun2025stylometric, bitton2025fingerprints, rosenfeld2026whose,
zaitsu2025stylometry, russell2026storyscope}. \citet{sun2025stylometric}
report 97.1\% over five chat models and 88.9\% after shuffling the words of
every text, so much of the signal lies in the word distribution rather than
its order. Attribution systems also concentrate false attributions on some
candidates \citep{alipoormolabashi2025maui}, as LLM judges do on GPT and
Claude. Few-shot LLM judges can approach classifier accuracy
\citep{dietrich2026splits}; we study the zero-shot case, which is how
anonymised judging pipelines use them.

\paragraph{Watermarks.}
Keyed text watermarks bias token choice using a secret key and the preceding
context \citep{kirchenbauer2023watermark, dathathri2024synthid}; detecting
them requires the key. Google marks text from the Gemini app with
SynthID-Text \citep{dathathri2024synthid}, and Anthropic announced keyed
watermarking of Claude text in August 2026. Section~\ref{sec:watermark}
explains why neither can account for our results.

\paragraph{Explanations.}
Judges' stated reasons often omit the cues that actually drove their
decisions \citep{marioriyad2025silent}. We analyse the 1,900 rationales our
frontier judges wrote with that caveat (Section~\ref{sec:rationales}).
